\section{Sim-to-Real Transfer}

\subsection{Por qué se necesita la simulación}

Hacer RL en robots reales plantea grandes desafíos:
\begin{itemize}
    \item La recolección de datos es lenta (se ejecuta en tiempo real)
    \item Riesgos de seguridad (el robot puede dañarse a sí mismo o dañar el entorno)
    \item Es difícil paralelizar a gran escala
\end{itemize}

Los entornos de simulación resuelven estos problemas, pero introducen la \textbf{sim-to-real gap}: la diferencia entre la simulación y el mundo real.

\begin{importantbox}{Orígenes de la Sim-to-Real Gap}
\begin{itemize}
    \item \textbf{Diferencias de dinámica}: el simulador no puede modelar a la perfección la física real (fricción, contacto, deformación)
    \item \textbf{Diferencias de percepción}: la visión y el tacto simulados difieren de los sensores reales
    \item \textbf{Diferencias de actuadores}: los motores reales tienen retardo, ruido y características no lineales
    \item \textbf{Diferencias de entorno}: el entorno real tiene objetos no modelados, cambios de iluminación, etc.
\end{itemize}
\end{importantbox}

\begin{figure}[H]
\centering
\includegraphics[width=0.9\textwidth]{slides-images/slide-12.jpg}
\caption{El ponente explica con un marco de entrenamiento con codificador de parámetros extrínsecos que aleatorizar en simulación factores como la masa, la fricción y el terreno es un paso de preparación esencial antes de transferir la política a un robot real.}
\end{figure}

\subsection{Domain Randomization}

\begin{importantbox}{Aleatorización de dominio}
Domain Randomization es uno de los métodos principales para cerrar la sim-to-real gap:

En la simulación se aleatorizan diversos parámetros del entorno (coeficientes de fricción, tamaño de los objetos, masa, apariencia visual, ruido de los sensores, etc.), de modo que la política, ya en la fase de entrenamiento, «haya visto» múltiples variantes posibles del entorno.

Supuesto central: si la política se desempeña bien en entornos simulados suficientemente diversos, el mundo real puede considerarse una de esas variantes.
\end{importantbox}

\begin{warningbox}{Limitaciones de Domain Randomization}
\begin{itemize}
    \item Una aleatorización excesiva puede volver la política demasiado conservadora
    \item Es necesario asegurar que los parámetros del mundo real estén dentro del rango de aleatorización
    \item Algunos fenómenos físicos (como el contacto con cuerpos blandos) son difíciles de cubrir incluso con aleatorización
    \item Elegir qué parámetros aleatorizar y en qué rango requiere conocimiento del dominio
\end{itemize}
\end{warningbox}

\subsection{Resumen de la sección}

La sim-to-real transfer es el puente fundamental para aplicar RL a robots reales, y la domain randomization es actualmente el método más utilizado.

